\documentclass[11pt]{article}
\usepackage[letterpaper,margin=1.03in]{geometry}
\usepackage{amsmath,amssymb,amsthm,mathtools,booktabs,array}
\usepackage[T1]{fontenc}
\usepackage{lmodern,microtype}
\usepackage{xcolor}
\usepackage[colorlinks=true,linkcolor=blue!45!black,urlcolor=blue!45!black,citecolor=blue!45!black]{hyperref}
\usepackage{fancyhdr}
\pagestyle{fancy}\fancyhf{}\fancyhead[L]{\small Filtered prime races}\fancyhead[R]{\small Research note / 26 November 2026}\fancyfoot[C]{\thepage}\setlength{\headheight}{14pt}
\newtheorem{theorem}{Theorem}
\newtheorem{lemma}[theorem]{Lemma}
\newtheorem{corollary}[theorem]{Corollary}
\newtheorem{proposition}[theorem]{Proposition}
\newcommand{\R}{\mathbb R}\newcommand{\C}{\mathbb C}\newcommand{\Z}{\mathbb Z}\newcommand{\e}{\mathrm e}\newcommand{\dd}{\,\mathrm d}\newcommand{\id}{\mathrm I}\newcommand{\ord}{\operatorname{ord}}\newcommand{\Rea}{\operatorname{Re}}\newcommand{\conv}{\operatorname{conv}}\newcommand{\norm}[1]{\lVert#1\rVert}\newcommand{\ph}{\varphi(q)}
\title{\bfseries Filtered prime races after GRH:\\zero-divisor rigidity and uniform multidirectional oscillation}
\author{Research note}
\date{26 November 2026}
\begin{document}
\maketitle
\begin{abstract}
We construct a compactly supported logarithmic averaging filter for the very ramp sum used in the recent Lean proof of the generalized Riemann hypothesis (GRH) for Dirichlet $L$-functions. The filter removes the pole main terms and any central-zero contribution, but retains every nonzero zero frequency. We prove, independently of GRH, that no nonzero complex linear combination of the zero multiplicity divisors of distinct primitive Dirichlet $L$-functions can have finite support. The proof uses Hadamard factorization and a two-unit shift of the logarithmic derivative. Under Dirichlet GRH, the resulting filtered prime-power count vector has a positive definite limiting second-moment matrix, even after any finite deletion of zero frequencies. Consequently every real unit contrast takes both positive and negative values of a uniform size on a uniform positive fraction of every sufficiently late logarithmic interval of a fixed length. The convex hull of the vector values on each such interval contains a fixed ball. No simplicity, disjointness, or rational independence of zeros is assumed. The filter is computed from finite prime sums alone. We provide a proof, a source audit, and a reproducible numerical illustration; the specific combination is proposed as a research contribution, not a certified priority claim.
\end{abstract}

\section{Source analysis and the choice of problem}
The starting point is the formal proof repository \cite{sernau}, inspected at commit
\texttt{72fead102f663d5c7}\allowbreak\texttt{a407964a13dc0cf9d977207}. Its top-level README describes a readable edition and a compact edition; the latter has about 76,000 lines of proof code. The top-level theorems in \path{grh-compact/src/Main.lean} prove Mathlib's RH statement, right-half-plane nonvanishing, and the critical-strip GRH statements for all Dirichlet characters. Their statement-only interface is \path{src/Statements.lean}. The README records the standard axioms and machine checks. The Lean FRO and Mathlib Initiative have publicly certified the original RH proof \cite{leanverification}; independent verification information is also collected by Rubinstein \cite{rubinstein}. This note uses the stated Dirichlet GRH theorem, not a new proof of it. We have inspected selected files, not rebuilt or independently kernel-checked the repository. The results here do not extend GRH to arbitrary automorphic or Dedekind zeta functions.

The most relevant reusable mechanism is the analytic endgame of the compact edition. In \path{New/Endgame/PrimeError.lean}, \texttt{lambdaClass} weights prime powers by $\Lambda(n)(1-\chi_K(n))$ in the class $1$ modulo $M$. The lemma \texttt{primeError\_eq\_sum} rewrites a normalized error as a finite sum with a ramp kernel, and \texttt{isPolar\_primeError} gives its Mellin/Laplace residues as differences of zero multiplicities. In \path{New/Endgame/OneSidedCriterion.lean}, the lemma \path{entireL_ne_zero_of_primeError_le} converts upper bounds into zero exclusion, using quadratic twists; \path{IsPolar.exists_window_sub_mul_le} locates a good second scale in a long window. \path{Lib/Analytic/Landau/PolarLandau.lean} supplies the one-sided transform machinery. In \path{New/Main.lean}, the arithmetic main inequality and the one-sided criterion yield \texttt{New.grh}.

The RH overview \cite{wiki} discusses error terms in prime distribution, analytic criteria and the geometry of the zeros. We reverse the direction of the endgame: knowing the critical line, what can the poles force about prime sums? An obstacle is cancellation between characters when zeros coincide. Rather than assuming distinct or linearly independent zeros, we show that complete cancellation of a tail is impossible. We then exploit compact logarithmic averaging to remove bias without evaluating $L'/L(1,\chi)$ or knowing the central-zero orders.

\section{Main statements}
Fix an integer $q\ge 1$ and let $G=(\Z/q\Z)^\times$ (the one-element group for $q=1$). The coordinates of all real vectors below are indexed by $a\in G$. Define the ramp sum, with the usual von Mangoldt function, by
\begin{equation}\label{eq:ramp}
 S_a(x)=\sum_{\substack{1\le n\le x\\n\equiv a\bmod q}}\Lambda(n)\left(\frac{x}{n}-1\right),\qquad
 R_a(t)=\e^{-t/2}S_a(\e^t),\quad t\in\R.
\end{equation}
The sum is zero for $x<1$ and the endpoint term is zero at $x=n$. Thus $R_a$ is continuous. Let $h>0$ be any fixed real parameter and define the averaging operator and scalar
\begin{equation}\label{eq:average}
 (A_h f)(t)=\frac1h\int_0^h f(t-u)\dd u,\qquad
 d_h=\frac{2(1-\e^{-h/2})}{h}\in(0,1).
\end{equation}
Define $D(t)=(D_a(t))_{a\in G}$ by the filter
\begin{equation}\label{eq:filter}
 \boxed{D_a(t)=(\id-A_h)(d_h\id-A_h)^2 R_a(t).}
\end{equation}
Only prime powers up to $\e^t$ enter $D(t)$. The filter uses $R_a$ over $[t-3h,t]$ and has no inputs from zero computations or $L$-values.

\begin{theorem}[Uniform prime-race oscillation]\label{thm:main}
Assume RH and GRH for Dirichlet $L$-functions. For every $q\ge 1$ and $h>0$, there exist $\epsilon,\eta,L_0>0$ and $T_0\in\R$, depending only on $q,h$, such that for every real unit vector $w\in\R^G$, every $T\ge T_0$, and every $L\ge L_0$,
\begin{align}
 \left|\{t\in[T,T+L]:w\cdot D(t)\ge\epsilon\}\right|&\ge\eta L,\label{eq:positive}\\
 \left|\{t\in[T,T+L]:w\cdot D(t)\le-\epsilon\}\right|&\ge\eta L.\label{eq:negative}
\end{align}
Here $|\cdot|$ is Lebesgue measure. The same constants work simultaneously for all directions $w$ (the times witnessing the inequalities may depend on $w$). In particular, with $K=\e^{L_0}$, each sufficiently late interval $[X,KX]$ contains both signs for every nonzero real contrast of the vector $D(\log x)$.
\end{theorem}

\begin{corollary}[A uniformly full-dimensional race body]\label{cor:convex}
For every $T\ge T_0$ and $L\ge L_0$,
\[
 \conv\{D(t):T\le t\le T+L\}\supseteq \{y\in\R^G:\norm{y}_2\le\epsilon\}.
\]
In particular no nontrivial affine hyperplane contains all the filtered vectors on such an interval. This is a convex-hull statement, not a claim that the trajectory itself visits every sector or every possible ordering.
\end{corollary}

\begin{theorem}[Zero-divisor tail rigidity, unconditional]\label{thm:rigidity}
For each character $\chi$ modulo $q$, let $\chi^*$ be its primitive inducing character. For a point $\rho$ in the open critical strip let $m_\chi(\rho)=\ord_\rho L(s,\chi^*)$ (zero if nonvanishing). If complex scalars $c_\chi$ satisfy
\[
 \sum_{\chi\bmod q}c_\chi m_\chi(\rho)=0
\]
for all but finitely many distinct nontrivial zero locations of the family, then every $c_\chi=0$. Equivalently, the multiplicity vectors $(m_\chi(\rho))_\chi$ outside any finite set span $\C^{\ph}$. This result does not assume GRH.
\end{theorem}

\section{Proof of zero-divisor rigidity}
We recall the standard analytic facts used here \cite{davenport,mv}: the primitive completed Dirichlet functions are entire of order one after removing the poles of the completed zeta function, their zeros are precisely the nontrivial zeros, and their zero counting functions satisfy $N_\chi(U)=O_q(U\log(U+2))$. Also the number in $U\le |\Im\rho|\le U+1$ is $O_q(\log(U+2))$. We use the standard meromorphic $L$-function, not a totalized value at a pole.

For $\chi^*$ of conductor $f_\chi$, let $a_\chi\in\{0,1\}$ be its parity and $\delta_\chi=1$ for the principal character and $0$ otherwise. Put
\[
 \Xi_\chi(s)=[s(s-1)]^{\delta_\chi}
 \left(\frac{f_\chi}{\pi}\right)^{s/2}
 \Gamma\left(\frac{s+a_\chi}{2}\right)L(s,\chi^*).
\]
Nonzero scalar factors are immaterial. The functional equation and nonvanishing on $\Rea s=1$ show that this completion has no extra zeros at $0$ or $1$. Its genus-one Hadamard expansion gives, locally normally off its zeros,
\begin{equation}\label{eq:hadamard}
 \frac{\Xi_\chi'}{\Xi_\chi}(s)=B_\chi+
 \sum_{\rho}m_\chi(\rho)\left(\frac1{s-\rho}+\frac1\rho\right).
\end{equation}
The convergence follows from $\sum_\rho m_\chi(\rho)|\rho|^{-2}<\infty$. All sums here and below distinguish a zero's location from its multiplicity.

\begin{proof}[Proof of Theorem \ref{thm:rigidity}]
Suppose the combined multiplicity is supported on a finite set. By \eqref{eq:hadamard}, the meromorphic function
\[
 B(s)=\sum_\chi c_\chi\frac{\Xi_\chi'}{\Xi_\chi}(s)
\]
is rational (a constant plus finitely many simple fractions). Define $F(s)=\sum_\chi c_\chi L'/L(s,\chi^*)$ and write $\psi=\Gamma'/\Gamma$. We have
\[
 B(s)=F(s)+\sum_\chi c_\chi\left[
 \tfrac12\log(f_\chi/\pi)+\tfrac12\psi((s+a_\chi)/2)
 +\delta_\chi(1/s+1/(s-1))\right].
\]
Now shift by \emph{two}. The identity $\psi(z+1)-\psi(z)=1/z$ means the gamma-factor contribution to $B(s+2)-B(s)$ is $\sum_\chi c_\chi/(s+a_\chi)$, a rational function. The remaining contributions are rational or cancel. Therefore
\begin{equation}\label{eq:rational}
 F(s+2)-F(s)\quad\text{is rational in }s.
\end{equation}
On the real axis as $\sigma\to+\infty$, the absolutely convergent Euler-product expansion
\begin{equation}\label{eq:euler}
 F(s)=-\sum_{n\ge2}\Lambda(n)\left(\sum_\chi c_\chi\chi^*(n)\right)n^{-s},\qquad\Rea s>1,
\end{equation}
shows $F(\sigma)=O_c(2^{-\sigma})$: for $\sigma\ge2$, bound $n^{-\sigma}\le 2^{2-\sigma}n^{-2}$ and sum. The rational function in \eqref{eq:rational} is thus exponentially small on the positive real axis. A nonzero rational function has leading behavior $C\sigma^k$ for some integer $k$ and $C\ne0$, so it must vanish identically. Hence $F(s+2)=F(s)$, and for each $s$ with $\Rea s>1$, letting the shift tend to infinity in \eqref{eq:euler} gives $F(s)=0$.

Uniqueness of absolutely convergent Dirichlet series gives $\sum_\chi c_\chi\chi^*(p)=0$ for every prime $p$ (the factor $\Lambda(p)$ is nonzero). For each reduced residue class $a$ modulo $q$ there is a prime $p\equiv a\bmod q$ by Dirichlet's theorem. Then $\chi^*(p)=\chi(a)$, so $\sum_\chi c_\chi\chi(a)=0$ for all $a$. Character orthogonality implies $c_\chi=0$ for every $\chi$. The spanning formulation is exactly the absence of a nonzero linear functional annihilating all those vectors.
\end{proof}

\paragraph{Why the argument matters.} Pairwise distinctness of some zeros would not by itself establish full rank of the entire family. The preceding proof excludes \emph{arbitrary complex linear cancellations} of multiplicities, even if zeros are shared, repeated or related. It also handles different conductors and both parities at once. The shift of length two is chosen to match both gamma factors.

\section{An explicit formula for the repository's ramp sum}
Write $S_\chi(x)=\sum_{n\le x}\Lambda(n)\chi(n)(x/n-1)$. Let
\[
 b_\chi=\lim_{s\to1}\left(-\frac{L'}{L}(s,\chi)-\frac{\delta_\chi}{s-1}\right)-\delta_\chi.
\]
The limit exists by the standard nonvanishing theorem at $1$ and the simple principal pole. The following absolutely convergent formula is useful for the filter.

\begin{lemma}[Ramp formula]\label{lem:ramp}
For $x\ge2$,
\begin{equation}\label{eq:explicit}
 S_\chi(x)=\delta_\chi x\log x+b_\chi x-
 \sum_\rho \frac{m_\chi(\rho)x^\rho}{\rho(\rho-1)}+O_q(1+\log x).
\end{equation}
The sum runs over distinct nontrivial zeros of the primitive inducing function, and the error is uniform in $x$ for fixed $q$. This formula itself is unconditional.
\end{lemma}
\begin{proof}
First suppose $\chi$ is primitive. Mellin inversion for the ramp kernel gives for $c=2$
\[
 S_\chi(x)=\frac{1}{2\pi i}\int_{c-i\infty}^{c+i\infty}
 -\frac{L'}{L}(s,\chi)\frac{x^s}{s(s-1)}\dd s.
\]
Indeed, the inverse Mellin integral of $y^s/[s(s-1)]$ on this line is $(y-1)_+$, including at $y=1$ by continuity. Absolute convergence of the Dirichlet series and the quadratic decay of the kernel justify exchange of sum and integral.

Shift to $\Rea s=-1/2$. The functional equation, Stirling's estimates and the absolutely convergent logarithmic derivative at $\Rea(1-s)=3/2$ show that on this line $L'/L(s,\chi)=O_q(\log(|\Im s|+2))$. There are no zeros or poles there. For completeness, the horizontal segments can be taken at heights $\pm U_j\to\pm\infty$ separated from every ordinate of the finitely many primitive functions by $\gg_q 1/\log U_j$. Such heights exist by the local unit-interval zero bound. On the horizontal segments $-1/2\le\Rea s\le2$, one has $L'/L(s,\chi)=O_q(\log^2 U_j)$. This usual local partial-fraction bound follows by subtracting \eqref{eq:hadamard} at $2+i\Im s$: zeros within distance one contribute $O_q(\log^2 U_j)$ by separation, while the remaining differences are $O_q(\sum_{k\ge1}\log(U_j+k+2)/k^2)=O_q(\log U_j)$; gamma and rational terms are bounded by Stirling. The horizontal integrals vanish for each fixed $x$ because the denominator is of order $U_j^2$. The new vertical integral is $O_q(x^{-1/2})$.

At each nontrivial zero $\rho$, the residue is $-m_\chi(\rho)x^\rho/[\rho(\rho-1)]$. At $1$, expand $-L'/L=\delta_\chi/(s-1)+b_\chi+\delta_\chi+O(s-1)$ and $x^s/s=x[1+(s-1)(\log x-1)+\cdots]$. The residue is $\delta_\chi x\log x+b_\chi x$. At $0$ there is at most a double pole (from the kernel and a possible trivial zero), contributing $O_q(1+\log x)$. No negative trivial zeros are crossed. Absolute convergence of the nontrivial-zero sum for fixed $x$ follows from the zero counting estimate and $0<\Rea\rho<1$.

For an imprimitive $\chi$ remove the Euler factors at $p\mid q$ from $\chi^*$. The exact difference is
\[
 S_{\chi^*}(x)-S_\chi(x)=
 \sum_{p\mid q}\log p\sum_{p^k\le x}\chi^*(p)^k(x/p^k-1)
 =x E_\chi+O_q(1+\log x),
\]
where $E_\chi=\sum_{p\mid q}\log p\sum_{k\ge1}\chi^*(p)^k p^{-k}$; terms at primes dividing the conductor vanish. To see the error, the omitted geometric tail multiplied by $x$ is bounded for each fixed $p$, and the subtracted $1$ terms contribute at most $\log x$. The Euler-factor identity gives $b_\chi=b_{\chi^*}-E_\chi$, and there is no change to the pole indicator. Substitution proves \eqref{eq:explicit}.
\end{proof}

Assume now Dirichlet GRH and write $\rho_\gamma=1/2+i\gamma$; let $\mathcal G\subset\R$ denote the union of distinct ordinates in the primitive family, including $0$ if it occurs. Put
\[
 u_\gamma(a)=-\frac{1}{\ph\rho_\gamma(\rho_\gamma-1)}
 \sum_\chi\overline{\chi(a)}\,m_\chi(\rho_\gamma).
\]
Orthogonality and \eqref{eq:explicit} yield, componentwise,
\begin{equation}\label{eq:R-expansion}
 R_a(t)=\frac{t\e^{t/2}}{\ph}+C_a\e^{t/2}
 +\sum_{\gamma\in\mathcal G}u_\gamma(a)\e^{i\gamma t}+O_q((1+t)\e^{-t/2}),
\end{equation}
where $C_a=\ph^{-1}\sum_\chi\overline{\chi(a)}b_\chi$. The constants are real by conjugation. The zero series converges uniformly and absolutely for all real $t$, since $|\rho_\gamma(\rho_\gamma-1)|=1/4+\gamma^2$. Also $u_{-\gamma}=\overline{u_\gamma}$, including for complex characters, by $m_{\overline\chi}(\rho_{-\gamma})=m_\chi(\rho_\gamma)$. The error estimate is used for large positive $t$ only.

\section{A filter that preserves the entire noncentral spectrum}
For $z\in\C$ set
\[
 a_h(z)=\frac{1-\e^{-hz}}{hz},\quad a_h(0)=1,\qquad
 M_h(z)=(1-a_h(z))(d_h-a_h(z))^2.
\]
The operator in \eqref{eq:filter} sends $\e^{zt}$ to $M_h(z)\e^{zt}$. Its multiplier has $M_h(0)=M_h(1/2)=M_h'(1/2)=0$. Differentiating the exponential identity in $z$ (the integration interval is compact) shows it annihilates $t\e^{t/2}$ as well as $\e^{t/2}$ and $1$. Thus both pole terms of \eqref{eq:R-expansion} and the central-zero term vanish. This is the key difference from merely subtracting a conjectural or estimated mean.

\begin{lemma}[No lost noncentral frequency]\label{lem:nonzero}
For every $h>0$ and real $\gamma\ne0$, $M_h(i\gamma)\ne0$.
\end{lemma}
\begin{proof}
The average of $\e^{-i\gamma u}$ over $[0,h]$ has modulus strictly smaller than one, because equality in the triangle inequality would require constant phase almost everywhere. Therefore $a_h(i\gamma)\ne1$. If $a_h(i\gamma)$ is real, its imaginary part $-(1-\cos(\gamma h))/(\gamma h)$ vanishes, so $\cos(\gamma h)=1$ and $\sin(\gamma h)=0$. In that case $a_h(i\gamma)=0$. Since $d_h>0$, $a_h(i\gamma)\ne d_h$. Both factors are nonzero.
\end{proof}

The filter has operator norm at most $2(1+d_h)^2$ on bounded functions and integrates only over a fixed length $3h$. Applying it to the expansion (termwise by uniform absolute convergence) proves
\begin{equation}\label{eq:D}
 D(t)=V(t)+O_{q,h}((1+t)\e^{-t/2}),\quad
 V(t)=\sum_{\substack{\gamma\in\mathcal G\\\gamma\ne0}}v_\gamma\e^{i\gamma t},\quad
 v_\gamma=M_h(i\gamma)u_\gamma.
\end{equation}
The error is in Euclidean norm. $V$ is a bounded, continuous, real vector-valued uniformly almost periodic function (a uniform limit of trigonometric polynomials), with $\sum_\gamma\norm{v_\gamma}_2<\infty$. The series of vectors includes multiplicities inside $v_\gamma$, not as independent terms with the same frequency.

\section{Positive definite second moments and uniform signs}
\begin{proposition}[Tail second moments]\label{prop:covariance}
For every $H\ge0$, define
\[
 V_H(t)=\sum_{|\gamma|>H}v_\gamma\e^{i\gamma t},\qquad
 Q_H=\sum_{|\gamma|>H}v_\gamma v_\gamma^*,
\]
where $^*$ denotes conjugate transpose. Under GRH, $Q_H$ is real symmetric positive definite. Uniformly in the starting point $T\in\R$,
\[
 \frac1L\int_T^{T+L}V_H(t)\dd t\longrightarrow0,\quad
 \frac1L\int_T^{T+L}V_H(t)V_H(t)^\top\dd t\longrightarrow Q_H
 \qquad(L\to\infty).
\]
\end{proposition}
\begin{proof}
For fixed $\omega\ne0$, $|L^{-1}\int_T^{T+L}\e^{i\omega t}\dd t|\le2/(L|\omega|)$, independently of $T$. This proves the limits for finite trigonometric polynomials, the second one by retaining only pairs $\gamma,-\gamma$. Uniform absolute convergence bounds the tail in supremum norm and hence the mean and product errors, proving the limits. Equivalently $Q_H$ is the Haar second moment on the closure of the frequency orbit; this formulation requires no independence assumption.

For $w\in\R^G$, $w^\top Q_Hw=\sum_{|\gamma|>H}|w\cdot v_\gamma|^2$. If this is zero then, by Lemma \ref{lem:nonzero},
\[
  \sum_\chi\left(\sum_{a\in G}w_a\overline{\chi(a)}\right)m_\chi(\rho_\gamma)=0
  \qquad(|\gamma|>H).
\]
Only finitely many locations are omitted. Theorem \ref{thm:rigidity} implies that every parenthesized Fourier coefficient is zero, and character inversion gives $w=0$. Thus the real symmetric matrix is positive definite. Conjugation pairs $\gamma,-\gamma$ justify reality. The same argument actually proves positive definiteness as a Hermitian matrix on complex vectors.
\end{proof}

\begin{proof}[Proof of Theorem \ref{thm:main}]
Use $V=V_0$, $Q=Q_0$. Set $B=\sum_{\gamma\ne0}\norm{v_\gamma}_2>0$ and $\lambda=\lambda_{\min}(Q)>0$. By Proposition \ref{prop:covariance}, choose $L_0$ so that for all $L\ge L_0$, all $T\in\R$, and every unit $w$, the scalar $g(t)=w\cdot V(t)$ satisfies
\[
 |g(t)|\le B,\qquad
 \left|\frac1L\int_T^{T+L}g\right|\le\frac{\lambda}{4B},\qquad
 \frac1L\int_T^{T+L}g^2\ge\frac{\lambda}{2}.
\]
Uniformity in $w$ follows from the vector norm and matrix operator norm convergence in the finite-dimensional space. Write $\mathbb E$ for averaging over the interval and $g_+=\max(g,0)$, $g_-=\max(-g,0)$. Since $g^2\le B|g|$,
\[
 \mathbb E g_\pm=\frac{\mathbb E|g|\pm\mathbb E g}{2}\ge\frac{\lambda}{8B}.
\]
Put $\theta=\lambda/(16B)$. The inequality $\mathbb E g_+\le\theta+B\,\mathbb P(g\ge\theta)$ gives $\mathbb P(g\ge\theta)\ge\lambda/(16B^2)$, and the same holds for $g\le-\theta$. Choose
\begin{equation}\label{eq:constants}
 \epsilon=\frac{\lambda}{32B},\qquad \eta=\frac{\lambda}{16B^2}.
\end{equation}
Choose $T_0$ large enough that $\norm{D(t)-V(t)}_2\le\epsilon$ for every $t\ge T_0$ using \eqref{eq:D}. Both measure assertions now follow by $|w\cdot(D-V)|\le\epsilon$, since $\theta=2\epsilon$. No limit distribution assumption or zero independence was needed.
\end{proof}
\begin{proof}[Proof of Corollary \ref{cor:convex}]
The continuous image of a compact interval has compact convex hull in finite dimension. Its support function in every unit direction is at least $\epsilon$ by Theorem \ref{thm:main}. If a point of the radius-$\epsilon$ ball were outside the hull, strict separation would produce a unit direction with its scalar product greater than the hull's support function, contradicting Cauchy--Schwarz. This proves the inclusion.
\end{proof}

\begin{corollary}[Robustness after deletion of low frequencies]\label{cor:tail}
For any fixed $H\ge0$, Theorem \ref{thm:main} and Corollary \ref{cor:convex} remain valid, with constants depending additionally on $H$, for the vector obtained from $D(t)$ by subtracting $\sum_{0<|\gamma|\le H} v_\gamma\e^{i\gamma t}$. In fact any finite conjugation-symmetric set of noncentral frequencies may be subtracted.
\end{corollary}
\begin{proof}
Use $V_H$ and $Q_H$ instead of $V,Q$. For any finite deletion the same proof of Proposition \ref{prop:covariance} works, since Theorem \ref{thm:rigidity} allows any finite exceptional set. The decaying error in \eqref{eq:D} is unchanged.
\end{proof}
This stronger version uses zero information for the subtraction; the basic theorem does not. In particular, spectral rank cannot be explained by any finite initial segment of zeros. For a matrix of real contrasts of full row rank, the induced second-moment matrix is also positive definite, by congruence with $Q_H$.

\section{Computation using prime powers only}
Define $b(v)=\mathbf 1_{v\ge0}(\e^{v/2}-\e^{-v/2})$. Then
\[
 R_a(t)=\sum_{n\equiv a\bmod q}\frac{\Lambda(n)}{\sqrt n}\,b(t-\log n),
\]
with finitely many nonzero terms. Expand the filter as
\[
 d_h^2\id-(d_h^2+2d_h)A_h+(1+2d_h)A_h^2-A_h^3.
\]
For $j\ge1$, let $b_j(v)=0$ for $v<0$, and for $v\ge0$ put
\[
 b_j(v)=2^j\left(\e^{v/2}-\sum_{k=0}^{j-1}\frac{(v/2)^k}{k!}\right)
 -(-2)^j\left(\e^{-v/2}-\sum_{k=0}^{j-1}\frac{(-v/2)^k}{k!}\right).
\]
It is the $j$-fold primitive with zero initial conditions. Taking $b_0=b$, repeated use of the fundamental theorem of calculus gives
\[
 A_h^jb(v)=h^{-j}\sum_{k=0}^j(-1)^k\binom jk b_j(v-kh).
\]
So $D_a(t)$ is evaluated by replacing $b$ in the finite prime-power formula by the linear combination of these four functions. For $v\ge3h$, the filter acts on the two exponentials with their multipliers, giving the numerically stable tail $-M_h(-1/2)\e^{-v/2}$. This avoids cancellation of growing terms.

The accompanying Python script uses a sieve and the finite-kernel identity, checks that the analytic kernel agrees with independent quadratures to numerical tolerance, and computes examples at $h=1$, 281 equally spaced points $5\le t\le12$. It uses 15,040 prime powers up to 162,751. Table \ref{tab:experiment} reports the empirical uncentered second-moment eigenvalue minimum and, for $q>1$, the extreme difference between the first and last unit-class coordinates. These are floating-point illustrations, not bounds for $Q$, not exact zero certificates, and not substitutes for the proofs.
\begin{table}[ht]
 \centering\small
 \begin{tabular}{r r r r}
 \toprule
 $q$&Smallest sample eigenvalue&Minimum contrast&Maximum contrast\\
 \midrule
 1&$4.084\cdot10^{-5}$&---&---\\
 4&$4.671\cdot10^{-5}$&$-0.05596$&$0.07152$\\
 5&$2.595\cdot10^{-5}$&$-0.10824$&$0.10399$\\
 8&$2.334\cdot10^{-5}$&$-0.12791$&$0.13041$\\
 \bottomrule
 \end{tabular}
 \caption{Illustrative finite samples; no inferential or asymptotic certification.}\label{tab:experiment}
\end{table}

\section{Novelty, limitations, and further questions}
The ingredients of the proof --- smoothed explicit formulas, character orthogonality, Hadamard products, and mean-square estimates for almost periodic functions --- are classical. Literature on prime races \cite{rs,mn,devin} treats sign changes, limiting distributions and the much stronger problem of inclusivity. There is also considerable work on common and distinct zeros \cite{fujii,banks}. The candidate contribution here is the specific all-contrasts, fixed-window theorem for a finite-prime filter of the Alp\"oge--Sernau ramp, together with the tail-rigidity proof via a rational shift and the robustness under finite spectral deletion. A targeted literature search did not identify this exact formulation. This is not an exhaustive priority search, and we do not claim that no equivalent lemma or theorem has appeared before.

The logical status is explicit: Theorem \ref{thm:rigidity} and the ramp identity are unconditional; the oscillation and covariance conclusions use Dirichlet GRH, now supplied by the reported verified development. The present proofs are classical mathematical proofs, not Lean-checked additions. The full original proof and its axiom audit were not rerun here. Constants in \eqref{eq:constants} are expressed via zero data and are positive, but no certified numerical $L_0,T_0,\epsilon,\eta$ are computed. The argument does not imply simplicity, rational linear independence, prime gaps beyond known GRH bounds, or inclusivity of ordinary prime races.

Several natural next steps are worth pursuing:
\begin{enumerate}
 \item Obtain effective lower bounds on $\lambda_{\min}(Q_H)$, ideally uniform in $q$ or in the discarded height $H$. The span proof implies a finite witness set but gives no height bound. A quantitative version of the shift argument or suitable twisted Landau sums might yield one.
 \item Derive a certified $L_0$ by combining a finite set of frequency gaps and an explicit zero-counting tail. Near-coincident distinct ordinates must be treated rigorously; floating-point zero comparisons do not settle multiplicities.
 \item Determine which alternative compact filters retain all noncentral frequencies while minimizing the size of the error or the necessary memory window; the double root at $1/2$ is forced by the pole term if no arithmetic constants are subtracted.
 \item Formalize the finite-dimensional mean-to-sign lemma and the multiplier argument in Lean, then connect the explicit formula and Hadamard shift to the analytic library in the pinned development. This would provide a checked corollary of the new GRH theorem without reproving the arithmetic main inequality.
 \item Investigate whether additional structural information about the zero frequencies makes the trajectory, not merely its convex hull, meet every ordering cone. The present proof deliberately avoids making that stronger claim.
\end{enumerate}

\begin{thebibliography}{99}\small
\bibitem{sernau} L. Sernau, collaborators and Claude, \emph{GRH for Dirichlet L-functions, in Lean 4}, 25 November 2026. Pinned repository: \url{https://github.com/lsernau/grh-lean/tree/72fead102f663d5c7a407964a13dc0cf9d977207}. See also the source map accompanying this note.
\bibitem{leanverification} Lean FRO and Mathlib Initiative, \emph{Lean verification of the Riemann Hypothesis}, November 2026. \url{https://leanprover-community.github.io/blog/posts/riemann-hypothesis/}.
\bibitem{rubinstein} M. Rubinstein, \emph{Alp\"oge/Claude proof of the Generalized Riemann Hypothesis}, verification log, November 2026. \url{https://www.math.uwaterloo.ca/~mrubinst/alpoge_claude_grh/}.
\bibitem{wiki} Wikipedia contributors, \emph{Riemann hypothesis}, consulted 26 November 2026. \url{https://en.wikipedia.org/wiki/Riemann_hypothesis}.
\bibitem{davenport} H. Davenport, \emph{Multiplicative Number Theory}, 3rd ed., revised by H. L. Montgomery, Springer, 2000. Classical Dirichlet $L$-function theory, zero counts, and explicit formulas.
\bibitem{mv} H. L. Montgomery and R. C. Vaughan, \emph{Multiplicative Number Theory I: Classical Theory}, Cambridge University Press, 2007. Dirichlet $L$-functions and explicit formula machinery.
\bibitem{rs} M. Rubinstein and P. Sarnak, Chebyshev's bias, \emph{Experimental Mathematics} 3 (1994), 173--197.
\bibitem{mn} G. Martin and N. Ng, Inclusive prime number races, \emph{arXiv:1710.00088}. \url{https://arxiv.org/abs/1710.00088}.
\bibitem{devin} L. Devin, Limiting properties of the distribution of primes in an arbitrarily large number of residue classes, \emph{Canadian Mathematical Bulletin}, DOI: \href{https://doi.org/10.4153/S0008439520000089}{10.4153/S0008439520000089}.
\bibitem{fujii} A. Fujii, On the zeros of Dirichlet $L$-functions (V), \emph{Acta Arithmetica} 28 (1976). \url{https://matwbn.icm.edu.pl/ksiazki/aa/aa28/aa2847.pdf}.
\bibitem{banks} W. Banks and K. Loftus, Unshared zeros of Dirichlet $L$-functions, \emph{arXiv:2607.22930}, 2026. \url{https://arxiv.org/abs/2607.22930}. Different problem: values of linear combinations of $L$-functions at zeros of one fixed function, rather than cancellation of logarithmic-derivative residues.
\end{thebibliography}
\end{document}
